%%%%%%%%%%%%%%%%%%%%%%%%%%%%%%%%%%%%%%%%%
\ifdefined\DISSERTATION
    \subsubsection{Indirect Optimal Control Solution Procedure}
    \label{sec:appendix-qp-numerical}
    The indirect optimal control approach applies optimality conditions and then numerically solves the resulting equations.
    The optimality conditions depend on the active set of constraints.
    In the most straightforward case where simply one constraint exists, the constraint is active if and only if the unconstrained solution violates that constraint.
    For example, if only a maximum force constraint exists, the constrained solution has a force $\gls{sym-hat-F-p-constr}$ that is the minimum of the unconstrained force $\gls{sym-hat-F-p-unconstrained}$ and the maximum allowed force fundamental $\gls{sym-F-max}$:
    \begin{equation}
    \gls{sym-hat-F-p-constr} = \min\left (\gls{sym-hat-F-p-unconstrained}, \gls{sym-F-max}\right)
    \end{equation}
    The solver guesses a value of $\gls{sym-Gamma}$, calculates the resulting force $\gls{sym-hat-F-p-guess}$, and then uses $\gls{sym-hat-F-p-guess}-\gls{sym-hat-F-p-constr}$ as an error signal to adjust its next guess of $\gls{sym-Gamma}$, iterating until the error converges to zero.
    This solver-based optimal control performs numerical root-finding, not numerical optimization, and can happen in an inner loop of the drag-linearization solver iteration (where $\gls{sym-hat-F-p-unconstrained}$ stays constant each control iteration) or simultaneously with the drag-linearization iteration (where $\gls{sym-hat-F-p-unconstrained}$ and the system linearization are updated each control iteration).
    MDOcean implements the latter, which simplifies the setup by requiring only one solver.
    Comparing the efficiency of the two approaches is left for future work.

    When multiple constraints exist, the situation is more complex, since it is possible for a given constraint to be violated by the unconstrained solution but not active in the constrained solution, or for multiple constraints to be active simultaneously.
    In the special case where the unconstrained solution violates no more than one constraint, that constraint is active and the solver can proceed as described above.
    When the unconstrained solution violates multiple constraints, the solver must try all subsets of 1-2 of those constraint (s) as the active set, find which of these active sets produces the largest power while satisfying all constraints, and then set the error as the deviation of the guessed solution from the constrained solution for that active set.
    Since the active set cannot be calculated a-priori as in the simpler case, this approach ends up looking more like a numerical optimization than a numerical root-finding, although the optimization is over a finite set of active set possibilities instead of a continuous control domain as it would be in a direct optimal control method.
    The main advantage of the indirect approach is its ability to handle arbitrary nonlinear\footnote{Nonlinear meaning greater than second order in an optimization sense, not in a dynamic sense, because the dynamics are already quasi-linearized here} constraints, although for the quadratic constraints considered here, the analytical approach described in \Cref{sec:appendix-qp-solution} is more efficient.
\fi
%%%%%%%%%%%%%%%%%%%%%%%%%%%%%%%%%%%%%%%%%
\subsubsection{Analytical Quadratic Program for Constrained Power Maximization}
\label{sec:appendix-qp-solution}
This section describes the analytical solution process for the constrained optimal control problem \Cref{eq:opt-problem} introduced in \Cref{sec:optimal-control}.

%%%%%%%%%%%%%%%
\paragraph{Quadratic Constraint Coefficients}
\label{sec:appendix-qp-constraints-table}

The quadratic coefficients $\{\gls{sym-mathbf-Q-mu}, \gls{sym-vec-a-mu}, \gls{sym-b-mu}\}$ in \Cref{eq:opt-problem} for each constraint implemented in MDOcean are given in \Cref{tab:qp-constraints}.

\begin{table}[b!]
\centering
\caption{Quadratic constraint coefficients}
\label{tab:qp-constraints}
\begin{tabular}{lccc}
    Constraint & $\gls{sym-mathbf-Q-mu}$ & $\gls{sym-vec-a-mu}$ & $\gls{sym-b-mu}$ \\ \hline
    $|\gls{sym-hat-xi-f}| \leq \gls{sym-xi-f-max}$ &
        $\frac{1}{\gls{sym-omega}^2}
        [\gls{sym-mathbf-a-F-p-dot-xi-f-leftarrow-VI}]^*
        \gls{sym-mathbf-A}_{|\gls{sym-hat-q}|^2}
        [\gls{sym-mathbf-a-F-p-dot-xi-f-leftarrow-VI}]$ &
        $[0,~ \gls{sym-hat-xi-f-0}]^\\text{top}$ & $\gls{sym-xi-f-max} - |\gls{sym-hat-xi-f-0}|$ \\
    $|\gls{sym-hat-xi-s}| \leq \gls{sym-xi-s-max}$ &
        $\frac{1}{\gls{sym-omega}^2}
        [\gls{sym-mathbf-a-F-p-dot-xi-s-leftarrow-VI}]^*
        \gls{sym-mathbf-A}_{|\gls{sym-hat-q}|^2}
        [\gls{sym-mathbf-a-F-p-dot-xi-s-leftarrow-VI}]$ &
        $[0,~ \gls{sym-hat-xi-s-0}]^\\text{top}$ & $\gls{sym-xi-s-max} - |\gls{sym-hat-xi-s-0}|$ \\
    $|\gls{sym-hat-X-PTO}| \leq \gls{sym-X-PTO-max}$ &
        $\frac{1}{\gls{sym-omega}^2}
        [\gls{sym-mathbf-a-F-dot-X-leftarrow-VI}]^*
        \gls{sym-mathbf-A}_{|\gls{sym-hat-q}|^2}
        [\gls{sym-mathbf-a-F-dot-X-leftarrow-VI}]$ &
        $\gls{sym-vec-0}$ & $\gls{sym-X-PTO-max}^2$ \\
    $|\gls{sym-hat-tau}| \leq \gls{sym-tau-max}$ &
        $[\gls{sym-mathbf-a-tau-Omega-leftarrow-VI}]^*
        \gls{sym-mathbf-A}_{|\gls{sym-hat-e}|^2}
        [\gls{sym-mathbf-a-tau-Omega-leftarrow-VI}]$ &
        $\gls{sym-vec-0}$ & $\gls{sym-tau-max}^2$ \\
    $\gls{sym-P-tau-Omega} \leq \gls{sym-P-tau-Omega-max}$ &
        $[\gls{sym-mathbf-a-tau-Omega-leftarrow-VI}]^*
        \gls{sym-mathbf-A}_{\gls{sym-P}}
        [\gls{sym-mathbf-a-tau-Omega-leftarrow-VI}]$ &
        $\gls{sym-vec-0}$ & $\gls{sym-P-tau-Omega-max}$ \\
    $\gls{sym-P-VI} \geq 0$ &
        $-\gls{sym-mathbf-A}_{\gls{sym-P}}$ &
        $\gls{sym-vec-0}$ & $0$ \\
\end{tabular}
\end{table}

%%%%%%%%%%%%%%%
\paragraph{Reduced Dimension Formulation}
The linear equality constraint $\gls{sym-vec-c}^{\,*}\gls{sym-vec-x} = d$ reduces the optimization problem dimension from $\gls{sym-vec-x} \\text{in} \gls{sym-mathbb-C}^2$ to $\gls{sym-mathbb-C}^1$.
Specifically, the equality constraint requires that the solution be a particular solution plus a scalar multiple of $\gls{sym-vec-n}$, the null space vector of $\gls{sym-vec-c}^{\,*}$:
\begin{equation}\label{eq:x-p-plus-nullspace}
    \gls{sym-vec-x} = \underbrace{
        \frac{\gls{sym-hat-V-s-th}}{2\Re(\gls{sym-Z-s-th})}
        \begin{bmatrix} \ Z_{s,th}^* \\ 1 \end{bmatrix}
    }_{
        \textrm{particular solution } \gls{sym-vec-x-p}}
    -~ \gls{sym-Gamma} \cdot \frac{\gls{sym-hat-V-s-th}}{2\Re(\gls{sym-Z-s-th})}
    \underbrace{
        \begin{bmatrix} -\ Z_{s,th}^* \\ 1 \end{bmatrix}
    }_{
        \textrm{null space } \gls{sym-vec-n}}
\end{equation}
where $\gls{sym-Gamma}$ is a nondimensional complex scalar representing the remaining degree of freedom in the optimization after enforcing the equality constraint.
For the particular solution, we have chosen to use the optimal solution in the absence of inequality constraints, which is identical to the solution implied by the impedance-matching condition \Cref{eq:matched-load} and can be derived from the quadratic program as follows:
\begin{equation}\label{eq:x-p}
    \gls{sym-vec-x-p} = \frac{d \gls{sym-mathbf-A-P}\gls{sym-vec-c}}{\gls{sym-vec-c}^{\,*}\gls{sym-mathbf-A-P}\gls{sym-vec-c}}
\end{equation}
This selection for $\gls{sym-vec-x-p}$ means that the degree of freedom $\gls{sym-Gamma}$ represents the deviation of the solution from the impedance-matched solution.
Using $\gls{sym-hat-I-p}$ for brevity to denote the current in the particular solution, \Cref{eq:x-p-plus-nullspace} can be rewritten as:
\begin{equation}\label{eq:x-p-plus-nullspace-simplified}
\gls{sym-vec-x} = \gls{sym-vec-x-p} - \gls{sym-Gamma} \gls{sym-hat-I-p} \gls{sym-vec-n}
    %&= \begin{bmatrix}
    %    1+\Gamma &  0 \\
    %    0 & 1-\Gamma
    %\end{bmatrix}
    %\vec{x}_p
\end{equation}

Inequality constraints
$\gls{sym-vec-x}^{\,*}\gls{sym-mathbf-Q-mu}\gls{sym-vec-x}
+ 2\Re\!\left\{\gls{sym-vec-a-mu}^{\,*}\gls{sym-vec-x}\right\}
\le \gls{sym-b-mu}$
further reduce the feasible design space, and the presence of multiple inequality constraints introduces the possibility of \Cref{eq:opt-problem} having no solutions due to infeasibility.
While constrained quadratic programs like \Cref{eq:opt-problem} generally require numerical solutions, the low dimensionality allows an analytical solution with geometric intuition, even if the problem is non-convex (see \Cref{sec:appendix-qp-convexity} for a discussion on convexity).
First, we rewrite the objective and inequality constraints in terms of $\gls{sym-Gamma}$ by plugging \Cref{eq:x-p-plus-nullspace-simplified} into \Cref{eq:opt-problem}:
\begin{subequations}
\begin{align}
\gls{sym-P-avg-VI}
&= \frac{|\gls{sym-hat-V-s-th}|^2}{8\Re (\gls{sym-Z-s-th})}~ (1 - |\gls{sym-Gamma}|^2)
\label{eq:obj-gamma}
\\[6pt]
\gls{sym-A-q-mu}|\gls{sym-Gamma}|^2 +&
2\Re\!\left\{
    \gls{sym-B-q-mu}\gls{sym-Gamma}
\right\}+\gls{sym-C-q-mu}
\le
0
\quad \forall \gls{sym-mu}
\label{eq:ineq-gamma}
\end{align}
\end{subequations}
where quadratic coefficients $\gls{sym-A-q-mu}$, $\gls{sym-B-q-mu}$, and $\gls{sym-C-q-mu}$ from \Cref{eq:ineq-gamma} are defined as:
\begin{equation}
\begin{aligned}
    \gls{sym-A-q-mu} &= |\gls{sym-hat-I-p}|^{2},
    \gls{sym-vec-n}^{\,*}\gls{sym-mathbf-Q-mu}\gls{sym-vec-n}& &\in \gls{sym-mathbb-R} \\
    \gls{sym-B-q-mu} &= -\gls{sym-hat-I-p},
    \left (\gls{sym-vec-x-p}^{\,*}\gls{sym-mathbf-Q-mu} + \gls{sym-vec-a-mu}^{\,*}\right)\gls{sym-vec-n}& &\in \gls{sym-mathbb-C} \\
    \gls{sym-C-q-mu} &=
        \gls{sym-vec-x-p}^{\,*}\gls{sym-mathbf-Q-mu}\gls{sym-vec-x-p}
        +
        2\Re\!\left\{\gls{sym-vec-a-mu}^{\,*}\gls{sym-vec-x-p}\right\}
        -\gls{sym-b-mu}&
    &\in \gls{sym-mathbb-R}
\end{aligned}
\end{equation}
Since \Cref{eq:obj-gamma} shows that power decreases with increasing $|\gls{sym-Gamma}|$, the power maximization problem \Cref{eq:opt-problem} can be recast as minimizing the norm of $\gls{sym-Gamma}$:
\begin{subequations}\label{eq:opt-problem-gamma}
\begin{alignat}{2}
& \min_{\gls{sym-Gamma} \in \gls{sym-mathbb-C}} &  |\gls{sym-Gamma}|^2& \\[3pt]
& \text{s.t.} &
S_{\mu}|\gls{sym-Gamma} - \Gamma_{c,\mu}|^2 &\le S_{\mu} r_{\mu}^2,
\quad \forall \mu \label{eq:disc-constraint}\\
& \text{where} &
\Gamma_{c,\mu}  &\in \gls{sym-mathbb-C} \nonumber\\
& & r_{\mu}^2 &\in \gls{sym-mathbb-R} \nonumber \\
& & S_{\mu} &\in \left\{-1,1\right\} \nonumber
\end{alignat}
\end{subequations}
where the constraint \Cref{eq:disc-constraint} is derived from \Cref{eq:ineq-gamma} by completing the square.
More precisely, when $\gls{sym-A-q-mu} \neq 0$, the $\gls{sym-mu}$th inequality constraint defines a circle in the complex plane of $\gls{sym-Gamma}$ centered at $\gls{sym-Gamma-c-mu}$ with radius $\gls{sym-r-mu-2}$.
The constraint requires that $\gls{sym-Gamma}$ lie within this circle when $\gls{sym-A-q-mu} > 0$ ($\gls{sym-S}_{\gls{sym-mu}}=1$) or outside this circle when $\gls{sym-A-q-mu} < 0$ ($\gls{sym-S}_{\gls{sym-mu}}=-1$).
Parameters $\gls{sym-Gamma-c-mu}$, $\gls{sym-r-mu-2}$, and $\gls{sym-S}_{\gls{sym-mu}}$ are defined as:
\begin{subequations}\label{eq:gamma-disc-params}
\begin{align}
\gls{sym-Gamma-c-mu-2}
&= \frac{-\gls{sym-B-q-mu}^{\,*}}{\gls{sym-A-q-mu}} =
\frac{
    \gls{sym-hat-I-p}^{\,*}\,
    \gls{sym-vec-n}^{\,*}
    \left (\gls{sym-mathbf-Q-mu}\gls{sym-vec-x-p} + \gls{sym-vec-a-mu}\right)
}{
    |\gls{sym-hat-I-p}|^{2}\,
    \gls{sym-vec-n}^{\,*}\gls{sym-mathbf-Q-mu}\gls{sym-vec-n}
}
\label{eq:gamma-disc-center}
\\[6pt]
\gls{sym-r-mu}^2
&=
\frac{|\gls{sym-B-q-mu}|^2-\gls{sym-A-q-mu}\gls{sym-C-q-mu}}{\gls{sym-A-q-mu}^2} =
\frac{
\gls{sym-b-mu}
-
\gls{sym-vec-x-p}^{\,*}\gls{sym-mathbf-Q-mu}\gls{sym-vec-x-p}
-
2\Re\!\left\{\gls{sym-vec-a-mu}^{\,*}\gls{sym-vec-x-p}\right\}
}{
|\gls{sym-hat-I-p}|^{2}\,
\gls{sym-vec-n}^{\,*}\gls{sym-mathbf-Q-mu}\gls{sym-vec-n}
}
+
|\gls{sym-Gamma-c-mu-2}|^2
\label{eq:gamma-disc-radius}
\\
\gls{sym-S-mu} &= \textrm{sign} (\gls{sym-A-q-mu}) = \textrm{sign}\left (\gls{sym-vec-n}^{\,*}\gls{sym-mathbf-Q-mu}\gls{sym-vec-n}\right)
\label{eq:gamma-disc-sign}
\end{align}
\end{subequations}

%%%%%%%%%%%%%%%
\paragraph{Solution}
\Cref{eq:opt-problem-gamma} can be interpreted graphically as finding the point closest to the origin that is interior to all constraint circles for which $\gls{sym-S-mu}=1$ and exterior to all circles for which $\gls{sym-S-mu}=-1$.
The unconstrained optimal solution of \Cref{eq:opt-problem-gamma} is the origin ($\gls{sym-Gamma}=0$), and the constrained optimal solution will also be zero in the case where the origin is feasible (no active constraints), which occurs when $ \gls{sym-S-mu} \gls{sym-C-q-mu} \leq 0 \quad \forall \gls{sym-mu}$.
If the origin is \textrm{infeasible}, the constrained optimal solution will lie on the boundary of one or more of the circles, specifically either at an intersection of two or more circles or at any circle's point closest to the origin.
We analytically compute all such candidate points and select the one with the lowest $|\gls{sym-Gamma}|$ that satisfies all constraints.
For $\gls{sym-N}$ constraints, there are up to $\gls{sym-N}^2$ candidate points ($\gls{sym-N}^2-\gls{sym-N}$ at intersections and $\gls{sym-N}$ at closest points).
This leads to the following expression for the constrained optimal solution (valid when $\gls{sym-A-q-mu} \neq 0 ~\forall \gls{sym-mu}$):
\begin{equation}\label{eq:gamma-opt-quadratic}
\gls{sym-Gamma-opt} =
\begin{cases}
 0 & \textrm{if } \gls{sym-S-mu} \gls{sym-C-q-mu} \leq 0 \quad \forall \gls{sym-mu} \\[3pt]
 \arg\min\limits_{\gls{sym-Gamma} \in \gls{sym-mathcal-F} \bigcap \gls{sym-mathcal-C}} |\gls{sym-Gamma}|
 & \textrm{otherwise} \\[3pt]
 \textrm{infeasible} & \textrm{if } \gls{sym-mathcal-F} \bigcap \gls{sym-mathcal-C} = \emptyset
\end{cases}
\end{equation}
where $\gls{sym-mathcal-F} = \bigcap_{\gls{sym-mu}} \left\{\gls{sym-Gamma}: \gls{sym-S}_{\gls{sym-mu}}|\gls{sym-Gamma} - \gls{sym-Gamma-c-mu}| \le \gls{sym-S}_{\gls{sym-mu}} \gls{sym-r-mu-2} \right\}$ is the set of all feasible points and $\gls{sym-mathcal-C}$ is the set of candidate points. $\gls{sym-mathcal-C}$ is the union of $\gls{sym-mathcal-C-int}$, the candidates at circle intersections, and $\gls{sym-mathcal-C-cl}$, the candidates at the closest point on each circle to the origin:
\begin{equation}\label{eq:candidate-points}
\begin{aligned}
\gls{sym-mathcal-C-int} &= \left\{
    \gls{sym-Gamma-mu-nu}, \quad \forall \gls{sym-mu} \neq \gls{sym-nu}
\right\}.
\\
\gls{sym-mathcal-C-cl} &= \left\{
    (|\gls{sym-Gamma-c-mu-2}| - \gls{sym-r-mu})\exp (i\angle\,\gls{sym-Gamma-c-mu-2}), \quad \forall \gls{sym-mu}
\right\}.
\end{aligned}
\end{equation}
The intersection points of the $\gls{sym-mu}$th and $\gls{sym-nu}$th circles are given by:
\begin{equation}\label{eq:circle-intersection}
    \gls{sym-Gamma-mu-nu} =
    \gls{sym-Gamma-c-mu-2} + \frac{\gls{sym-Gamma-c-nu-2} - \gls{sym-Gamma-c-mu-2}}{\gls{sym-d-mu-nu}}
    \left(        \gls{sym-a-mu-nu}
        \pm
        1i \sqrt{\gls{sym-r-mu}^2-\gls{sym-a-mu-nu}^2}
    \right)
\end{equation}
where $\gls{sym-d-mu-nu-star} = |\gls{sym-Gamma-c-nu} - \gls{sym-Gamma-c-mu}|$ is the distance between the circle centers,
$\gls{sym-a-mu-nu-star} = \frac{1}{2}(\gls{sym-r-mu-2}^2-\gls{sym-r-nu}^2+\gls{sym-d-mu-nu-star}^2)/\gls{sym-d-mu-nu-star}$ is the distance from
the center of the $\gls{sym-mu}$th circle to the midpoint of the two intersection points,
and $1i$ is the imaginary unit.
The closest-point candidates $\gls{sym-mathcal-C-cl}$ match the solution that \citet[\equationautorefname~4]{mccabe_force-limited_2024} derive for problems with only one constraint.

Importantly, we need not compute the entire feasible space $\gls{sym-mathcal-F}$, only compute $\gls{sym-mathcal-C}$ and then check whether each candidate point is feasible.
\Cref{fig:qp-solution-geometry} illustrates the various cases for feasibility and optimality for a toy problem with $\gls{sym-r-mu-2}=4$, $\gls{sym-S}_{\gls{sym-mu}}=1$, and various $\gls{sym-Gamma-c-mu}$.
\begin{figure}[htbp]
    \centering
    \includegraphics[\FigureOptions{aor:figs/from-matlab/qp_circles.pdf}]{figs/from-matlab/qp_circles.pdf}
    \caption{Visualization of the optimization problem \Cref{eq:opt-problem-gamma} in the complex plane of $\gls{sym-Gamma}$.
Toy problems with various constraint circle centers $\gls{sym-Gamma-c-mu}$, all with radius $\gls{sym-r-mu-2}=4$ and $\gls{sym-S}_{\gls{sym-mu}}=1$, demonstrate the various solution cases in \Cref{eq:gamma-opt-quadratic}.}
    \label{fig:qp-solution-geometry}
\end{figure}

If $\gls{sym-A-q-mu} = 0$ for any $\gls{sym-mu}$, the $\gls{sym-mu}$th inequality constraint defines a line in the complex plane of $\gls{sym-Gamma}$ rather than a circle.
This can occur when $\gls{sym-mathbf-Q}_{\gls{sym-mu}} = 0$ (the constraint on $\gls{sym-vec-x}$ is linear rather than quadratic, which is not the case for any of the constraints in \Cref{tab:qp-constraints}) or when $\gls{sym-vec-n}^{\, *}\gls{sym-mathbf-Q}_{\gls{sym-mu}}\gls{sym-vec-n} = 0$ ($\gls{sym-mathbf-Q}_{\gls{sym-mu}}$ is rank-deficient and $\gls{sym-vec-n}$ is in its nullspace).
The same types of critical points (now consisting of line-line, circle-line, and circle-circle intersections, as well as the closest point on each circle and line) apply, but for implementation simplicity, we replace $\gls{sym-A-q-mu}$ with a small positive $\gls{sym-epsilon}$ in this case and use the quadratic implementation.

%%%%%%%%%%%%%%%
\paragraph{Quartic Apparent Power Constraints}
\label{sec:appendix-mag-S-constraints}
While constraints on effort and flow magnitudes and real and reactive power are quadratic in $\gls{sym-vec-x}$, constraints on the magnitude of apparent power $|\gls{sym-S}|<|\gls{sym-S}_{\text{max}}|$, and therefore the maximum and minimum instantaneous power, are fourth-order in $\gls{sym-vec-x}$ and thus $\gls{sym-Gamma}$.
This results in a ``Cassini oval'' curve \citep{weisstein_cassini_2026} on the complex plane of $\gls{sym-Gamma}$, which consists of two closed curves when $|\gls{sym-S}_{\text{max}}|<|\gls{sym-S-0}|$ and one curve when $|\gls{sym-S}_{\text{max}}|\geq|\gls{sym-S-0}|$, where $|\gls{sym-S-0}|$ is the apparent power of the unconstrained optimum.
Though the optimization is no longer a quadratically constrained quadratic program, the same geometric solution procedure applies, adding circle-oval intersections $\gls{sym-mathcal-C}_{int,\text{oval}-\text{circ}}$, oval-oval intersections, and oval closest points $\gls{sym-mathcal-C}_{cl,\text{oval}}$ to the set of candidate points $\gls{sym-mathcal-C}$.
The new candidate points are:
\begin{equation}
\begin{aligned}
\gls{sym-mathcal-C-int-oval-circ} &= \left\{
    \gls{sym-Gamma}:
    |\gls{sym-Gamma} - \gls{sym-Gamma-c-mu-2}| = \gls{sym-r-mu}, ~(1-|\gls{sym-Gamma}|^2)^2 + 4\Im (\gls{sym-Gamma}) = \frac{|\gls{sym-S-max-k}|}{|\gls{sym-S-0}|}, \quad \forall \gls{sym-mu},\gls{sym-k-generic-index}
\right\}.
\\
\gls{sym-mathcal-C-cl-oval} &= \left\{
    \gls{sym-Gamma}:
    \gls{sym-Gamma} = \begin{cases}
        \pm \sqrt{1-\frac{|\gls{sym-S-max-k}|}{|\gls{sym-S-0}|}}, \quad \forall \gls{sym-k-generic-index}~ |\gls{sym-S-max-k}|<|\gls{sym-S-0}| \\
        \pm \sqrt{1+\frac{|\gls{sym-S-max-k}|}{|\gls{sym-S-0}|}}, \quad \forall \gls{sym-k-generic-index}~ |\gls{sym-S-max-k}|>|\gls{sym-S-0}|
    \end{cases}
\right\}.
\end{aligned}
\end{equation}
where $\gls{sym-mu}$ indexes circle constraints and $\gls{sym-k-generic-index}$ indexes oval constraints.
While the earlier circle-circle intersections of \Cref{eq:candidate-points} result in a quadratic equation that is solved analytically for the coordinates of $\gls{sym-Gamma}$, oval-circle intersections lead to quartic equations that are solved numerically with standard polynomial root-finding routines.
Typically there is only one apparent power constraint, so oval-oval intersections are not derived, but if present would likewise yield a fourth-order polynomial in $\gls{sym-Gamma}$.
Therefore, the presence of apparent power constraints makes the optimal control method semi-analytical rather than analytical.

%%%%%%%%%%%%%%%
\paragraph{Extension to Higher Dimensions and Convexity}
\label{sec:appendix-qp-convexity}
Future work may wish to extend the constrained power maximization to the full nonlinear dynamics, or to linear dynamics with irregular wave forcing as in the study by \citet{bacelli_numerical_2015}.
In the terminology of \Cref{fig:venn-diagram}, these would be known as nonlinear and linear spectral optimal control respectively.
Both of these cases solve the full harmonic-balance equation, producing a higher-dimensional analog of \Cref{eq:opt-problem-gamma}.
The problem dimension equals twice the number of frequencies retained in the spectral expansion, preventing the use of the geometric analytical solution used in the monochromatic case.
Instead, a numerical solver must be used, and the computational scaling of that solver depends critically on whether the problem is \emph{convex}.

A convex optimization problem is one in which both the feasible set and the objective function are convex---informally, the feasible region has no holes or concavities, and the objective is bowl-shaped with no local valleys.
The key practical implication is that any locally optimal solution is also globally optimal, so gradient-based solvers converge to the true optimum without exhaustive search, and solution time scales polynomially with the number of decision variables rather than exponentially.

After reducing to the $\gls{sym-Gamma}$ parameterization, the power objective $|\gls{sym-Gamma}|^2$ is a sum of squares and therefore convex.
Even though the original power matrix $\gls{sym-mathbf-A-P}$ (see \Cref{tab:power-metrics}) has eigenvalues $\{-1,1\}$ and is indefinite, \citet{bacelli_numerical_2015} show that substituting the linear dynamics equality constraint into the objective---which is exactly what the reduction to $\gls{sym-Gamma}$ in \Cref{eq:obj-gamma} accomplishes---renders the objective convex for both regular and irregular wave forcing.
Nonlinear dynamics can reintroduce non-convexity and should be checked on a case-by-case basis.

The inequality constraints are more problematic.
A constraint is convex if the set of points satisfying it forms a convex region.
In the $\gls{sym-Gamma}$ plane, a convex constraint corresponds to a circle whose feasible region is the \emph{interior} ($\gls{sym-S-mu}=1$ in \Cref{eq:disc-constraint}), while a non-convex constraint corresponds to a circle whose feasible region is the \emph{exterior} ($\gls{sym-S-mu}=-1$).
Constraints on the effort and flow magnitude (e.g., force limits $|\gls{sym-hat-F}|\leq \gls{sym-F}_{\text{max}}$) can be convex or non-convex, depending on the complex angle of the Th\'{e}venin equivalent of the relevant source impedance and on the ratio of the constraint limit to its unconstrained value.
\citet{mccabe_force-limited_2024} visualize the feasible regions for families of such constraints on the complex $\gls{sym-Gamma}$ plane.
They show that the constraints are always convex for purely real impedances; that for reactive plants, nonconvexities occur when the value of the constraint limit exceeds some threshold; and that this threshold lowers as the reactivity increases.
The quadratic matrices $\gls{sym-mathbf-Q}$ in \Cref{tab:qp-constraints} confirm this: $\gls{sym-mathbf-A}_{|\gls{sym-hat-q}|^2}$ and $\gls{sym-mathbf-A}_{|\gls{sym-hat-e}|^2}$ are only positive \emph{semi}definite, indicating constraints that define half-spaces rather than bounded ellipsoids.
For the 2D problem studied here, non-convexity is handled by the exhaustive enumeration of candidate points in \Cref{eq:gamma-opt-quadratic}, which is efficient in two dimensions but does not scale to higher dimensions.
Future work should investigate semidefinite relaxations of these non-convex constraints using the S-procedure \citep{vanantwerp_tutorial_2000}, which replaces each non-convex constraint with a conservative but convex outer approximation.
It is expected that the optimal design will have low reactivity because reactivity increases signal peak-to-average ratios, implying that a semi-definite relaxation would likely improve in accuracy as it approaches the optimal.

%%%%%%%%%%%%%%%
\paragraph{Comparison to Other WEC Constrained Control Formulations}
Notably, $\gls{sym-Gamma}$ is equal to the reflection coefficient at the generator electrical port, which is commonly used in microwave circuit design to quantify impedance mismatch.
\citet{mccabe_force-limited_2024,coe_co-design_2025} further discuss the reflection coefficient for WECs.

Similar graphical interpretations of constraints for wave energy converters are presented by \citet{mccabe_force-limited_2024,merigaud_geometrical_2023,bacelli_geometric_2013}, with all three visualizing quadratic constraints as circles on the 2D plane.
\citet{mccabe_force-limited_2024} focus on force/current limits but discuss amplitude, phase voltage, and apparent power constraints as well.
\citet{merigaud_geometrical_2023} consider exclusively amplitude limits while \citet{bacelli_geometric_2013} cover amplitude and force limits.
\citet{mccabe_force-limited_2024} use Smith charts to show constraints directly in $\gls{sym-Gamma}$ space as we do here, while \citet{merigaud_geometrical_2023} work in the complex plane of the hydrodynamic transmission coefficient and \citet{bacelli_geometric_2013} in that of the PTO force.
\citet{bacelli_geometric_2013} also extend the formulation to capture higher harmonics (nonsinusoidal forces), where the exact constraint cannot be easily represented geometrically but sufficient conditions for constraint violation and satisfaction in the higher dimensional space can be represented as hyperspheres and hyperellipsoids.

%%%%%%%%%%%%%%%%%%%%%%%%%%%%%%%%%%%%%%%%%

\subsubsection{Constraint Sensitivity Analysis via Parametric Programming}
\label{sec:appendix-constraint-sensitivity}
To support the discussion in \Cref{sec:discussion}, this section derives the scaling of optimal average power with the constraint limits and demonstrates its conditions for convexity.

The constraint limits are
$\gls{sym-tau}_{\text{max}}$, $\gls{sym-P}_{\gls{sym-tau}\gls{sym-Omega},\text{max}}$, $\gls{sym-xi}_{\gls{sym-f},\text{max}}$, $\gls{sym-xi}_{\gls{sym-s},\text{max}}$, and $\gls{sym-X-PTO-max}$, which we collectively denote as $\gls{sym-L-mu}$.
The dependence of power on the first two limits (PTO sizing variables) is of more immediate relevance than the last three (amplitude limits) because the former can be independently manipulated through PTO design, while the latter can only be manipulated by changing the bulk dimensions.
The sensitivity of power to bulk dimensions is substantially more complex than the sensitivity to constraint limits for reasons discussed in \Cref{sec:appendix-dimension-sensitivity}, so we focus on the PTO sizing sensitivities here.

In a given sea state, the global sensitivity of optimal power to a constraint limit can be quantified using parametric programming theory, where a local sensitivity is found parametrically and then integrated over the parameter space to find the global parameter relationship.
The local sensitivity of power to a constraint parameter
$\frac{\partial \gls{sym-p}_{\text{avg},VI,\text{opt}}}{\partial \gls{sym-L-mu}}$
equals the product of the Lagrange multiplier for that constraint, $\gls{sym-lambda-mu}$,
and the partial derivative of the constraint function with respect to the parameter,
$\frac{\partial \gls{sym-g-mu}}{\partial \gls{sym-L-mu}}$ \cite[Section~5.9.3]{boyd_convex_2004}.
Since the constraint function is defined as
$\gls{sym-g-mu} (\gls{sym-vec-x}) = \gls{sym-vec-x}^{\,*}\gls{sym-mathbf-Q-mu}\gls{sym-vec-x}
                 + 2\Re\!\left\{\gls{sym-vec-a-mu}^{\,*}\gls{sym-vec-x}\right\} - \gls{sym-b-mu}\leq 0$,
and \Cref{tab:qp-constraints} shows that each constraint limit $\gls{sym-L-mu}$ enters $\gls{sym-g-mu}$ only through $\gls{sym-b-mu}$, we have
$\frac{\partial \gls{sym-g-mu}}{\partial \gls{sym-L-mu}} = -\frac{\partial \gls{sym-b-mu}}{\partial \gls{sym-L-mu}}$.
This derivative can be trivially computed for all constraints in \Cref{tab:qp-constraints}.
The power sensitivity simplifies to:
\begin{equation}
    \frac{\partial \gls{sym-p-avg-VI-opt}}{\partial \gls{sym-L-mu}} = -\gls{sym-lambda-mu} \frac{\partial \gls{sym-b-mu}}{\partial \gls{sym-L-mu}}
\end{equation}
If the optimization problem were solved with a numerical solver,
the Lagrange multipliers would be directly available as part of the solution.
In our case, we have a semi-analytical solution for the optimal reflection coefficient at each sea state in terms of the constraint centers and radii at that sea state,
$\gls{sym-Gamma-opt}(\gls{sym-r-mu-2}, \gls{sym-Gamma-c-mu})$,
so we can compute the Lagrange multipliers as follows:
\begin{equation}\label{eq:lagrange}
\begin{aligned}
   \gls{sym-lambda-mu}
    &\equiv
    \frac{\partial \gls{sym-p-avg-VI-opt}}{\partial \gls{sym-g-mu}}
    \\
    &= - \frac{\partial \gls{sym-p-avg-VI-opt}}{\partial \gls{sym-b-mu}}
    \\
    &= -
    \frac{\partial \gls{sym-p-avg-VI-opt}}{\partial |\gls{sym-Gamma-opt}|}
    \frac{\partial |\gls{sym-Gamma-opt}|}{\partial \gls{sym-b-mu}}
    \\
    &= -
    \frac{\partial \gls{sym-p-avg-VI-opt}}{\partial |\gls{sym-Gamma-opt}|}
    \left(        \frac{\partial |\gls{sym-Gamma-opt}|}{\partial \gls{sym-r-mu}}
        \frac{\partial \gls{sym-r-mu}}{\partial \gls{sym-b-mu}}
        +
        \frac{\partial |\gls{sym-Gamma-opt}|}{\partial \gls{sym-Gamma-c-mu-2}}
        \frac{\partial \gls{sym-Gamma-c-mu-2}}{\partial \gls{sym-b-mu}}
    \right)
\end{aligned}
\end{equation}
The computation of each of the five partials in the final expression of \Cref{eq:lagrange} is discussed next.

%%%%%%%%%%%%%%%
\paragraph{Effect of Constraint Limits on Circle Radii}
Recall from \Cref{eq:gamma-disc-radius} that each inequality constraint maps to a circle in the complex $\gls{sym-Gamma}$ plane
with center $\gls{sym-Gamma-c-mu}$ and radius $\gls{sym-r-mu}$ given by:
\begin{equation}\label{eq:radius-full}
    \gls{sym-r-mu}^2
    = \frac{\gls{sym-b-mu} - \gls{sym-vec-x-p}^{\,*}\gls{sym-mathbf-Q-mu}\gls{sym-vec-x-p}
            - 2\Re\!\left\{\gls{sym-vec-a-mu}^{\,*}\gls{sym-vec-x-p}\right\}}
           {|\gls{sym-hat-I-p}|^{2}\,\gls{sym-vec-n}^{\,*}\gls{sym-mathbf-Q-mu}\gls{sym-vec-n}}
    + |\gls{sym-Gamma-c-mu-2}|^2
\end{equation}
The constraint limit $\gls{sym-L-mu}$ enters $\gls{sym-r-mu}$ only through $\gls{sym-b-mu}$
 (the right-hand side of the quadratic inequality in \Cref{tab:qp-constraints}),
while $\gls{sym-vec-x-p}$, $\gls{sym-vec-n}$, $\gls{sym-hat-I-p}$, $\gls{sym-mathbf-Q-mu}$, $\gls{sym-vec-a-mu}$, and $\gls{sym-Gamma-c-mu}$ depend only on the Th\'{e}venin impedance $\gls{sym-Z-s-th}$ and excitation $\gls{sym-hat-V-s-th}$ for the sea state,
which are independent of the constraint limits.
Defining $\gls{sym-d-mu} = |\gls{sym-hat-I-p}|^{2}\,\gls{sym-vec-n}^{\,*}\gls{sym-mathbf-Q}_{\gls{sym-mu}}\gls{sym-vec-n}$
and $\gls{sym-c-mu} = \gls{sym-d-mu}|\gls{sym-Gamma-c-mu}|^2 - \gls{sym-vec-x-p}^{\,*}\gls{sym-mathbf-Q}_{\gls{sym-mu}}\gls{sym-vec-x-p} - 2\Re\!\left\{\gls{sym-vec-a-mu-2}^{\,*}\gls{sym-vec-x-p}\right\}$,
both independent of the constraint limits, \Cref{eq:radius-full} and its derivative simplify to:
\begin{equation}\label{eq:radius-simplified}
    \gls{sym-r-mu}^2 = \frac{\gls{sym-b-mu} + \gls{sym-c-mu}}{\gls{sym-d-mu}}
, \quad
    \frac{\partial \gls{sym-r-mu}}{\partial \gls{sym-b-mu}} = \frac{1}{2\sqrt{\gls{sym-d-mu} (\gls{sym-b-mu} + \gls{sym-c-mu})}}
\end{equation}
Since $\gls{sym-r-mu}\in\gls{sym-mathbb-R}$ (circle radii cannot have an imaginary part), then $\gls{sym-b-mu} + \gls{sym-c-mu}$ and $\gls{sym-d-mu}$ must have the same \textrm{sign}, so $\frac{\partial \gls{sym-r-mu}}{\partial \gls{sym-b-mu}}$ is always positive.
In other words, increasing $\gls{sym-b-mu}$ always enlarges the circle radius, regardless of whether the constraint represents the interior or exterior of the circle.
Meanwhile, \Cref{eq:gamma-disc-center} shows that the center $\gls{sym-Gamma-c-mu}$ does not depend on $\gls{sym-b-mu}$, so $\frac{\partial \gls{sym-Gamma-c-mu}}{\partial \gls{sym-b-mu}}=0$.
Thus, the constraint limit $\gls{sym-L-mu}$ does not affect the circle center.

%%%%%%%%%%%%%%%
\paragraph{Optimal Reflection Coefficient as a Function of Constraint Limits}
Now we combine the relationship between circle radius and optimal reflection coefficient $\gls{sym-Gamma-opt}(\gls{sym-r-mu})$, with the relationship between circle radius and constraint limit $\gls{sym-r-mu}(\gls{sym-b-mu})$ derived above
to obtain a relationship between reflection coefficient magnitude and constraint limit in a given sea state, $|\gls{sym-Gamma-opt}|(\gls{sym-b-mu})$.
Substituting \Cref{eq:radius-simplified} into \Cref{eq:candidate-points} shows that
the squared optimal reflection coefficient magnitude in each sea state takes one of the following values depending on the number of active constraints:
\begin{equation}\label{eq:gamma-opt-summary}
    |\gls{sym-Gamma-opt}|^2 =
    \begin{cases}
        0
    & \text{none active} \\
        |\gls{sym-Gamma-c-mu}|^2
        - 2|\gls{sym-Gamma-c-mu}|\sqrt{\frac{\gls{sym-b-mu-2}+\gls{sym-c-mu-2}}{\gls{sym-d-mu-2}}}
        + \frac{\gls{sym-b-mu-2}+\gls{sym-c-mu-2}}{\gls{sym-d-mu-2}}
    & \text{only $\gls{sym-mu}^*$ active} \\
        \left|
            \Gamma_{c,\mu^*} + \frac{\Gamma_{c,\nu^*} - \Gamma_{c,\mu^*}}{d_{\mu^*\nu^*}}
            \left(                \gls{sym-a-mu-nu-star}
                \pm
                1i \sqrt{\frac{\gls{sym-b-mu-2} + \gls{sym-c-mu-2}}{\gls{sym-d-mu-2}}-\gls{sym-a-mu-nu-star}^2}
            \right)
        \right|^2
    & \text{$\gls{sym-mu}^*$ \& $\gls{sym-nu}^*$ active}
    \end{cases}
\end{equation}
where $\gls{sym-mu}^*,\gls{sym-nu}^*$ refer to the indices of the active constraint (s) for a given sea state,
and $\gls{sym-a}_{\gls{sym-mu}^*,\gls{sym-nu}^*}$ and $d_{\gls{sym-mu}^*,\gls{sym-nu}^*}$ are real numbers defined in the discussion following \Cref{eq:circle-intersection}.
$d_{\gls{sym-mu}^*,\gls{sym-nu}^*}$ is independent of the constraint limits, while $\gls{sym-a}_{\gls{sym-mu}^*,\gls{sym-nu}^*}$ is affine in $\gls{sym-b}_{\gls{sym-mu}^*}$ and $\gls{sym-b}_{\gls{sym-nu}^*}$.
Thus, using \textrm{aff}() and \textrm{quad}() to denote affine and quadratic functions of the arguments, we have the following scale behavior:
\begin{equation}\label{eq:gamma-opt-scale}
    |\gls{sym-Gamma-opt}|^2 =
    \begin{cases}
        0
    & \text{none active} \\
        \textrm{aff} (\gls{sym-b-mu-2}) + \sqrt{\textrm{aff} (\gls{sym-b-mu-2})}
    & \text{only $\gls{sym-mu}^*$ active} \\
        \left|
            \textrm{aff} (\gls{sym-b-mu-2})+\textrm{aff} (\gls{sym-b-nu-2})
            \pm
            1i \sqrt{\textrm{quad} (\gls{sym-b-mu-2}, \gls{sym-b-nu-2})}
        \right|^2
    & \text{$\gls{sym-mu}^*$ \& $\gls{sym-nu}^*$ active} \\
    = \textrm{quad} (\gls{sym-b-mu-2}, \gls{sym-b-nu-2}) +
    (\textrm{aff} (\gls{sym-b-mu-2}) + \textrm{aff} (\gls{sym-b-nu-2}))
    \sqrt{\textrm{quad} (\gls{sym-b-mu-2}, \gls{sym-b-nu-2})} &
    \end{cases}
\end{equation}
Constraint activity indices $\gls{sym-mu}^*,\gls{sym-nu}^*$ depend on the constraint limits,
and this dependence must be accounted for before the curvature of $|\gls{sym-Gamma-opt}|^2$ can be analyzed.
From \Cref{eq:gamma-opt-quadratic}, we can define $\gls{sym-mu}^*$ and $\gls{sym-nu}^*$ more explicitly as functions of $\gls{sym-b-mu}$ by defining a penalty function $\gls{sym-f}(\gls{sym-Gamma})$ that returns $\gls{sym-Gamma}$ if $\gls{sym-Gamma}$ is feasible and $\infty$ \textrm{otherwise}.
This function can be constructed as $\gls{sym-f}(\gls{sym-Gamma}) = \gls{sym-Gamma} + \sum_{\gls{sym-mu}} \gls{sym-mathbb-I-mu}(\gls{sym-Gamma})$,
where $\gls{sym-mathbb-I-mu}(\gls{sym-Gamma})$ is an indicator function that returns $0$ if $\gls{sym-Gamma}$ satisfies the $\gls{sym-mu}$th constraint and $\infty$ \textrm{otherwise}.
The penalty function is convex in $\gls{sym-Gamma}$ if and only if all constraints are convex, i.e., if $\gls{sym-S-mu}=1 ~\forall \gls{sym-mu}$.
Then, the optimal reflection coefficient of each candidate type can be written as the minimum of the candidate points passed through the penalty function.
\begin{equation}\label{eq:gamma-opt-min-per-candidate}
\begin{aligned}
   |\gls{sym-Gamma-cl-mu}| &=
    \min_\gls{sym-mu} \left\{
        | \gls{sym-f}\left (\gls{sym-Gamma-cl-mu-2} (\gls{sym-b-mu})\right) |
    \right\}, \quad
    |\Gamma_{int,\mu^*,\nu^*}| =
    \min_{\mu,\nu} \left\{
        | f\left (\Gamma_{\mu,\nu} (b_\mu, b_\nu)\right) |
    \right\}
\end{aligned}
\end{equation}
Finally, the overall optimal reflection coefficient $|\gls{sym-Gamma-opt}|$ is the minimum of all the per-candidate-type minima:
\begin{equation}\label{eq:gamma-opt-min}
\begin{aligned}
    |\gls{sym-Gamma-opt}| &=
    \min \left\{
        |\Gamma_{cl,\mu^*} (b_\mu)|,  |\Gamma_{int,\mu^*,\nu^*} (b_\mu, b_\nu)|
    \right\}
\end{aligned}
\end{equation}

Working from the inside out, we now analyze the convexity of $|\gls{sym-Gamma-opt}|^2$ with respect to $\gls{sym-b-mu}$.
Starting with the candidates at closest points, the square root is a concave increasing function,
and a concave increasing function of an affine function is concave,
so $\sqrt{\textrm{aff}(\gls{sym-b-mu-2})}$ is concave.
The sum of a concave function and an affine function is concave,
so $|\gls{sym-Gamma-cl}|=\textrm{aff}(\gls{sym-b-mu-2}) + \sqrt{\textrm{aff}(\gls{sym-b-mu-2})}$ is concave.
The penalty function is convex if and only if all constraints are convex,
so $-|\gls{sym-f}(\gls{sym-Gamma-cl})|$ is concave if and only if $\gls{sym-S-mu}=1 ~\forall \gls{sym-mu}$.
The minimum of a set of concave functions is a concave increasing function,
% or the maximum of a set of convex functions is a convex increasing function
and $|\gls{sym-Gamma}_{cl,\gls{sym-mu}^*}|$ is the minimum of a set of concave functions if and only if all constraints are concave,
so $|\gls{sym-Gamma}_{cl,\gls{sym-mu}^*}|$ is concave if and only if $\gls{sym-S-mu}=1 ~\forall \gls{sym-mu}$.

Moving onto the candidates at intersection points, a quadratic function is concave if its leading order term is negative, and the leading order term of $\textrm{quad}(\gls{sym-b-mu-2}, \gls{sym-b-nu-2})$ is negative (arising from the subtraction of squared-affine term $\gls{sym-a-mu-nu-star}^2$).
%Square root is a concave function, and a concave function of a concave function is concave, so $\sqrt{\textrm{quad}(\gls{sym-b-mu-2}, \gls{sym-b-nu-2})}$ is concave.
The geometric mean of two positive values is a concave function, so the geometric mean of two positive concave functions is concave.
The term $\left(\textrm{aff}(\gls{sym-b-mu-2}) + \textrm{aff}(\gls{sym-b-nu-2})\right) \sqrt{\textrm{quad}(\gls{sym-b-mu-2}, \gls{sym-b-nu-2})}$
is the geometric mean of $(\textrm{aff}(\gls{sym-b-mu-2})+\textrm{aff}(\gls{sym-b-nu-2}))^2$ and $\textrm{quad}(\gls{sym-b-mu-2}, \gls{sym-b-nu-2})$.
The $(\textrm{aff}(\gls{sym-b-mu-2})+\textrm{aff}(\gls{sym-b-nu-2}))^2$ is positive but is convex rather than concave.
This means the curvature of $\left(\textrm{aff}(\gls{sym-b-mu-2}) + \textrm{aff}(\gls{sym-b-nu-2})\right) \sqrt{\textrm{quad}(\gls{sym-b-mu-2}, \gls{sym-b-nu-2})}$ cannot be established with the present analysis.
If the curvature of $\left(\textrm{aff}(\gls{sym-b-mu-2}) + \textrm{aff}(\gls{sym-b-nu-2})\right) \sqrt{\textrm{quad}(\gls{sym-b-mu-2}, \gls{sym-b-nu-2})}$ could be established as concave,
then the sum of two concave functions is concave, so the candidate points $|\gls{sym-Gamma-int}|$ would be concave.
If that were the case, then by the same argument as for the closest points, the minimum of feasible $|\gls{sym-Gamma-int}|$ would be concave if and only if all constraints are concave.
The optimum $|\gls{sym-Gamma-opt}|$ is the minimum of the feasible $|\gls{sym-Gamma-cl}|$ and $|\gls{sym-Gamma-int}|$, so if both of those are concave, then it is also concave.
Therefore, the criteria for the concavity of $|\gls{sym-Gamma-opt}|$ with respect to $\gls{sym-b-mu}$ is that all constraints are convex (interior of circles), and the curvature of $\left(\textrm{aff}(\gls{sym-b-mu-2}) + \textrm{aff}(\gls{sym-b-nu-2})\right) \sqrt{\textrm{quad}(\gls{sym-b-mu-2}, \gls{sym-b-nu-2})}$ is concave.
As future work, the above curvature analysis should be verified with a disciplined convex programming package such as CVXPY.

To facilitate the analysis of the next section, we observe that although the squared optimal reflection coefficient magnitude, $|\gls{sym-Gamma-opt}|^2$,
has so far been expressed as a function only of the active constraints $(\gls{sym-mu}^*,\gls{sym-nu}^*)$,
it can be written as a sum over all constraint pairs $(\gls{sym-mu},\gls{sym-nu})$ multiplied by an appropriate Kronecker delta function:
\begin{equation}
\label{eq:gamma-opt-scale-sum}
\begin{aligned}
    |\gls{sym-Gamma-opt}|^2
    = &\sum_\gls{sym-mu} \sum_{\gls{sym-nu} > \gls{sym-mu}}
    \left[
        \delta_{\mu,\mu^*}(1-\delta_{\nu,\nu^*})
        \left(            \textrm{aff} (\gls{sym-b-mu})
            + \sqrt{\textrm{aff} (\gls{sym-b-mu})}
        \right)
     + \delta_{\mu,\mu^*}\delta_{\nu,\nu^*}
    \right.
\\
    &\left.
        \left(            \textrm{quad} (\gls{sym-b-mu}, \gls{sym-b-nu}) +
            \left (\textrm{aff} (\gls{sym-b-mu}) + \textrm{aff} (\gls{sym-b-nu})\right)
            \sqrt{\textrm{quad} (\gls{sym-b-mu}, \gls{sym-b-nu})}
        \right)
    \right]
\end{aligned}
\end{equation}
where $\gls{sym-delta}_{\gls{sym-mu},\gls{sym-mu}^*}\gls{sym-delta}_{\gls{sym-nu},\gls{sym-nu}^*}$ is non-zero when both $\gls{sym-mu}$ and $\gls{sym-nu}$ are active in that sea state,
and $\gls{sym-delta}_{\gls{sym-mu},\gls{sym-mu}^*}(1-\gls{sym-delta}_{\gls{sym-nu},\gls{sym-nu}^*})$ is non-zero when $\gls{sym-mu}$ is active and $\gls{sym-nu}$ is inactive in the sea state.
Both Kronecker delta functions are implicitly functions of the constraint limit $\gls{sym-b-mu}$, as seen in \Cref{eq:gamma-opt-min}.
% delta mu* is monotonially decreasing (negative step) with $\gls{sym-b-mu}$ when $\gls{sym-S-mu}=1$ (enlarging the mu-th circle when the mu-th constraint is feasible-inside will make more candidate points mu-feasible, so the mu-th constraint would transition from active to inactive)
% and monotonically increasing with $\gls{sym-b-mu}$ when $\gls{sym-S-mu}=-1$ (enlarging the mu-th circle when the mu-th constraint is feasible-outside will make fewer candidate points mu-feasible, so the mu-th constraint would transition from inactive to active).
% $\gls{sym-delta}_{\gls{sym-nu},\gls{sym-nu}^*}$ is monotonially decreasing with $\gls{sym-b-mu}$.

%%%%%%%%%%%%%%%
\paragraph{Average Power as a Function of Constraint Limits}
The results just obtained establish conditions for the convexity of power in a single sea state to the constraint limit in that sea state, which we now make explicit by adding a $\gls{sym-beta}$ subscript to all sea state dependent quantities.
For LCOE, we are interested in the weighted average power across sea states as a function of the constraint limits.
All constraint limits are independent of the sea state ($\gls{sym-L-mu-beta}=\gls{sym-L-mu}$), with the exception of the amplitude limits due to slamming.
\Cref{sec:appendix-slam} will show that the slamming limits are a function of $\gls{sym-Delta} \gls{sym-z}_{\text{slam}}/\gls{sym-H}$, which depends explicitly on the wave height, and $\gls{sym-theta}$, which depends implicitly on both the wave period and wave height through the guessed phase $\angle\gls{sym-xi}$.
Since the PTO sizing limits are of primary relevance, the following derivation assumes that the constraint limits are independent of the sea state to simplify the expression, although sea-state-dependent limits can easily be accommodated by adding a $\gls{sym-beta}$ index on $\gls{sym-b-mu}$ and carrying it through the derivation.

From \Cref{eq:power-elec-sum,eq:gamma-opt-summary}, the average power can be written as:
\begin{equation}\label{eq:power-elec-sensitivity}
    \overline{\gls{sym-P}}_{elec} = \overline{\gls{sym-P}}_{elec,0}
    - \gls{sym-eta} \sum_{\gls{sym-beta}=1}^{\gls{sym-N-sea}} \gls{sym-w-beta}
    \cdot |\gls{sym-Gamma-beta-opt}|^2
\end{equation}
where $\overline{\gls{sym-P}}_{\text{elec},0}$ is the power available if all sea states were unconstrained (impedance-matched) and
$\gls{sym-w-beta}=JPD_{\gls{sym-beta}} \cdot
    \frac{|\gls{sym-hat-V-s-th-beta}|^2}
         {8\Re (\gls{sym-Z-s-th-beta})}$ is a constant weight for each sea state.
The second term represents power lost to constraint-induced impedance mismatch and depends on the constraint limits through \Cref{eq:gamma-opt-summary}.

Substituting \Cref{eq:gamma-opt-scale-sum} into \Cref{eq:power-elec-sensitivity} requires a $\gls{sym-beta}$ subscript on the starred index to the Kronecker delta functions since constraint activity varies by sea state, as well as a $\gls{sym-beta}$ subscript on the affine and quadratic functions to capture sea-state dependence of the Th\'{e}venin parameters:
\begin{equation}
\begin{aligned}
    \overline{P}_{elec} = &\overline{P}_{elec,0}
    - \eta \sum_\mu \sum_{\nu > \mu}
    \sum_{\beta=1}^{N_{sea}} w_\beta
    \left[
        \delta_{\mu,\mu^*_\beta}(1-\delta_{\nu,\nu^*_\beta})
        \left(            \textrm{aff}_\beta (b_{\mu})
            + \sqrt{\textrm{aff}_\beta (b_{\mu})}
        \right)
    \right.
\\
    &\left.
        + \delta_{\mu,\mu^*_\beta}\delta_{\nu,\nu^*_\beta}
        \left(            \textrm{quad}_\beta (b_{\mu}, b_{\nu}) +
            \left (\textrm{aff}_\beta (b_{\mu}) + \textrm{aff}_\beta (b_{\nu})\right)
            \sqrt{\textrm{quad}_\beta (b_{\mu}, b_{\nu})}
        \right)
    \right]
\end{aligned}
\end{equation}
Expanding under the assumption that each constraint is active by itself in some sea states and active with another constraint in others,
we find that $\gls{sym-overline-P-elec}$ is a piecewise function of each constraint limit $\gls{sym-b-mu}$, with a quadratic leading order term as well as affine, square root of affine, and affine times square root of quadratic terms.
The two Kronecker delta terms are combined with the weighting term $\gls{sym-w-beta}$ to create revised weighting terms $\gls{sym-tilde-w-1-mu-nu-beta}$ and $\gls{sym-tilde-w-2-mu-nu-beta}$ that are non-zero only when the corresponding constraint (s) are active in that sea state:
\begin{equation}\label{eq:power-triple-sum}
\begin{aligned}
    \gls{sym-overline-P-elec} &= \overline{P}_{elec,0}
    - \eta \sum_\mu \sum_{\nu > \mu}
    \left[
        \sum_{\beta=1}^{N_{sea}} \tilde{w}_{1\mu\nu\beta}
        \left(            \textrm{aff}_\beta (b_{\mu})
            + \sqrt{\textrm{aff}_\beta (b_{\mu})}
        \right)
    \right. + \\
    &\left.
     \sum_{\beta=1}^{N_{sea}} \tilde{w}_{2\mu\nu\beta}
        \left(            \textrm{quad}_\beta (b_{\mu}, b_{\nu}) +
            \left (\textrm{aff}_\beta (b_{\mu}) + \textrm{aff}_\beta (b_{\nu})\right)
            \sqrt{\textrm{quad}_\beta (b_{\mu}, b_{\nu})}
        \right)
    \right]
\end{aligned}
\end{equation}
so the power-constraint relationship is a sum over sea states of piecewise functions of the constraint limit $\gls{sym-b-mu}$.
Recognizing that the sea state sum in \eqref{eq:power-triple-sum} essentially just scales the individual piecewise functions, the expression can be simplified by absorbing the weights into the subscripts of the affine and quadratic terms, and combining the sum of an affine and quadratic term into a single quadratic term:
\begin{equation}\label{eq:power-double-sum}
\begin{aligned}
    \gls{sym-overline-P-elec} = \overline{P}_{elec,0}
    - \eta \sum_{\mu\nu\beta}
    &\left[
        \sqrt{\textrm{aff}_{\mu\nu\beta} (b_{\mu})}
        + \textrm{quad}_{\mu\nu\beta} (b_{\mu}, b_{\nu})
    +
        \left (\textrm{aff}_{\mu\nu\beta} (b_{\mu}) + \textrm{aff}_{\mu\nu\beta} (b_{\nu})\right)
        \sqrt{\textrm{quad}_{\mu\nu\beta} (b_{\mu}, b_{\nu})}
    \right]
\end{aligned}\end{equation}
We can then use a similar analysis as before to obtain conditions for the convexity of the overall power-constraint relationship.

%The convexity of $|\gls{sym-Gamma-opt}|^2$ in $\gls{sym-b-mu}$ derived above translates directly to convexity of $\gls{sym-overline-P-elec}$ in $\gls{sym-b-mu}$ via \Cref{eq:power-elec-sensitivity}, since the weighted sum of concave functions ($-|\gls{sym-Gamma-opt}|^2$) is concave and the addition of constant $\overline{\gls{sym-P}}_{\text{elec},0}$ preserves concavity. Equivalently, $\gls{sym-overline-P-elec}$ is convex in the constraint coefficients $\gls{sym-b-mu}$ under the same conditions, namely that all constraints are convex ($\gls{sym-S-mu} = 1$) and that the curvature of $(\textrm{aff}(\gls{sym-b-mu}) + \textrm{aff}(\gls{sym-b-nu}))\sqrt{\textrm{quad}(\gls{sym-b-mu}, \gls{sym-b-nu})}$ resolves favorably (see the caveat above). This is the convexity claim referenced in \Cref{sec:multidisciplinary-insights}.

% \paragraph{Convexity of the Power-Constraint Relationship}
% To illustrate convexity, we use an intuitive geometric argument, leaving formal proof to future work.
% In a constraint that has $\gls{sym-S-mu}=1$ in a given sea state, the feasible region is the interior of a circle, so increasing the constraint limit $\gls{sym-b-mu}$ enlarges the circle radius and therefore enlarges the feasible region, which can only decrease or maintain the optimal $|\gls{sym-Gamma-beta-opt}|$ and therefore increase or maintain power.
% In a constraint with $\gls{sym-S-mu}=-1$, the feasible region is the exterior of a circle, so increasing the constraint limit $\gls{sym-b-mu}$ enlarges the circle radius and therefore shrinks the feasible region, which can only increase or maintain the optimal $|\gls{sym-Gamma-beta-opt}|$ and therefore decrease or maintain power.
% This corresonds to the \textrm{sign} of the leading order term of a given piecewise segment in the power-constraint relationship, which is negative for $\gls{sym-S-mu}=1$ and positive for $\gls{sym-S-mu}=-1$.
% Furthermore, the number of sea states in which a particular constraint is singularly or pairwise active ($\sum_{\gls{sym-beta}}\gls{sym-mathbb-I}_\gls{sym-beta}$) is monotonically non-decreasing with the constraint limit, so the leading order term in the overall power-constraint relationship (all piecewise segments together) is positive if the constraint has $\gls{sym-S-mu}=1$ at all sea states, or negative if $\gls{sym-S-mu}=-1$ at all sea states.
% This means that as long as the \textrm{sign} of $\gls{sym-S-mu}$ is consistent across sea states, the power-constraint relationship is globally convex.
% Relaxing any constraint thus always increases or maintains power, and the marginal power benefit diminishes as the constraint limit grows.

%%%%%%%%%%%%%%%
\paragraph{Extension to Dimension Sensitivity}
\label{sec:appendix-dimension-sensitivity}
In principle, the same parametric programming approach can be used to semi-analytically compute the derivatives of the constrained optimal power with respect to bulk dimensions,
though the implementation remains out of reach.
These derivatives would be highly valuable, and indeed would enable the full wave-to-wire control co-design optimization problem to be solved semi-analytically, rather than the present approach which requires a numerical bulk-dimension optimization to be wrapped around the semi-analytical control optimization and simulation.
Computing bulk-dimension derivatives would require not only $\partial \gls{sym-Gamma-opt}/\partial \gls{sym-b-mu}$ as we derive here,
but also the derivatives of Th\'{e}venin parameters with respect to bulk dimensions,
because bulk dimensions affect $\gls{sym-mathbf-Q-mu}, \gls{sym-vec-a-mu}$, and the unconstrained power through the Th\'{e}venin parameters.
The sensitivity formula would become significantly more complex and less intuitive with these additional terms,
and more importantly, the Th\'{e}venin-to-dimension derivatives are not currently available in MDOcean because they depend on the hydrodynamic coefficient-to-dimension derivatives.
Application of the adjoint method has recently enabled efficient computation of hydrodynamic coefficient derivatives in the context of a numerical boundary element method solver \citep{khanal_fully_2025},
but the adjoint method has not yet been implemented for the semi-analytical MEEM model used here.
Implementing the adjoint for a linear solve is relatively straightforward in a multidisciplinary design optimization package like OpenMDAO,
and would unlock a fully semi-analytical differentiation workflow.


%%%%%%%%%%%%%%%%%%%%%%%%%%%%%%%%%%%%%%%%%
\subsubsection{Force Saturation}
\label{sec:appendix-force-sat}
\Cref{sec:dynamics} introduces the use of describing functions to quantify the fundamental amplitude of non-sinusoidal force waveforms, which can achieve higher powers than the equivalent sinusoidal waveform for the same force limit.

The simulation permits the nonsinusoidal waveform to have a maximum fundamental a factor of $\frac{4}{\gls{sym-pi}}$ above the force limit.
This coincides with the fundamental to peak ratio of a square wave, which is the limiting case both for the bang-singular-bang controller (a piecewise-discontinuous sine-square wave combination that \citet{hendrikx_optimal_2017} finds as the analytically optimal solution), and for the saturated sine controller (a piecewise-continuous sine-square wave combination that \citet{coe_initial_2020} finds as numerically optimal).
In fact, the optimization need not specify the exact nonlinear waveform, since under the describing function filtering hypothesis, only the fundamental matters for simulation purposes.
For hardware implementation purposes, the nonlinear controller can be obtained for the saturated-sine controller by the process outlined in the study \cite{mccabe_force-limited_2024}.

The appropriate complex scale factor from the optimal unsaturated controller (damping or reactive respectively) is found as described in \Cref{sec:appendix-qp-solution} and will be real for the damping control case and complex in the reactive control case.

%Mathematically, a saturation factor $0< \gls{sym-f}_{sat}\leq 1$ is computed as the ratio between the fundamental amplitude of the saturated torque and the original unsaturated torque signal,
%\begin{equation}
%    f_{sat} = \min\left(\alpha\frac{\tau_{max}}{\tau}, 1\right),
%\end{equation}
%where $\gls{sym-tau}_{\text{max}}$ is the maximum permissible generator torque and $\gls{sym-alpha}$ is the ratio between the fundamental amplitude of the saturated waveform and the maximum torque $\gls{sym-tau}_{\text{max}}$, using the describing function derived in \citep{mccabe_force-limited_2024},
%\begin{equation}
%    \alpha = \frac{2}{\pi} \left(\frac{F_p}{\tau_{max}}\sin^{-1}\frac{\tau_{max}}{F_p} +  \sqrt{1 - \left (\frac{\tau_{max}}{F_p}\right)^2} \right).
%\end{equation}

%When $\frac{\gls{sym-F-p}}{\gls{sym-tau}_{\text{max}}}$ = 1 (no saturation), then $\gls{sym-alpha}= 1$;
%when $\frac{\gls{sym-F-p}}{\gls{sym-tau}_{\text{max}}} \rightarrow \infty$ (full saturation), then $\gls{sym-alpha}= \frac{4}{\gls{sym-pi}}$.
%This makes sense because $\frac{4}{\gls{sym-pi}}$ is the harmonic amplitude of a unit square wave,
%and in full saturation the saturated signal approaches a square wave.
%$\gls{sym-alpha}$ is plotted in \Cref{fig:alpha-fit}.

%The factor $\gls{sym-f-sat}$ describes the ratio of the saturated and unsaturated torque signal.
%Previous work \citet{mccabe_multidisciplinary_2022} implies that this is also the ratio of the saturated and unsaturated control gains, $\gls{sym-m-sat}$, but this is not the case because the WEC response also changes with saturation.

% %Substituting $\gls{sym-f}_{sat}$ back into the equation of motion reveals that the control gain saturation multiplier $\gls{sym-m-sat}$ is the solution to a quadratic equation $\gls{sym-a-q} \gls{sym-m-sat}^2 + \gls{sym-b-q} \gls{sym-m-sat} + \gls{sym-c-q} = 0$. If the uncontrolled dynamics have a natural frequency $\gls{sym-omega}_{\gls{sym-n},\gls{sym-u}}$ and damping ratio $\gls{sym-zeta-u}$, and the controlled unsaturated dynamics have natural frequency $\gls{sym-omega-n}$ and damping ratio $\gls{sym-zeta}$, then the quadratic coefficients can be expressed as:

% \begin{equation}\label{eq:abc-matrix}
%     \begin{bmatrix}
%         a_q \\ b_q \\ c_q
%     \end{bmatrix}
%     =
%     \begingroup % keep the row spacing change local
%     \setlength\arraycolsep{2pt}
%     \begin{bmatrix} -4 & \frac{4}{f_{sat}^2} & \frac{1}{f_{sat}^2} & -1 & 0 & 0\\
%     8 & 0 & 0 & 2 & 8 & 2\\
%     -4 & -4 & -1 & -1 & -8 & -2 \end{bmatrix}
%     \endgroup
%     \begin{bmatrix}
%  (\theta_1)^2  \\
%   (\theta_2)^2 \\
%    (\theta_3)^2 \\
%    (\theta_4)^2  \\
%    \theta_1 \theta_2 \\
%     \theta_3 \theta_4
%     \end{bmatrix}
% \end{equation}
% where the basis vector $\vec{\gls{sym-theta}}$ is a dimensionless function of the natural frequencies and damping ratios:
% \begin{equation}\label{eq:theta-basis}
%         \begin{bmatrix}
%   \theta_1  \\
%   \theta_2 \\
%    \theta_3 \\
%    \theta_4  \\
%     \end{bmatrix} =     \begin{bmatrix}
% ~ \overline{\omega} (\zeta_u - \overline{\omega}_n\zeta)  \\
%   \overline{\omega} ~\overline{\omega}_n \zeta \\
%   \overline{\omega}^2 - \overline{\omega}_n^2\\
%    \overline{\omega}_n^2 - 1  \\
%     \end{bmatrix}
% \end{equation}
% and the overlined frequencies indicate nondimensionalization by the uncontrolled natural frequency: $\overline{\gls{sym-omega}} = \frac{\gls{sym-omega}}{\gls{sym-omega}_{\gls{sym-n},\gls{sym-u}}}, \overline{\gls{sym-omega}}_{\gls{sym-n}} = \frac{\gls{sym-omega-n}}{\gls{sym-omega}_{\gls{sym-n},\gls{sym-u}}}$. The quadratic equation is then solved for $\gls{sym-m-sat}$, and of the two roots, the one that is real-valued and between 0 and 1 is used.

% The no saturation case ($\gls{sym-f}_{sat} = 1$) for \eqref{abc-matrix} yields $\gls{sym-m-sat}=1$. This confirms the expected behavior that when no force reduction is required, the powertrain coefficients can remain unchanged.

% Equations \eqref{eq:abc-matrix} and \eqref{eq:theta-basis} are an equivalent reformulation of \citet{mccabe_force-limited_2024}, which uses expressions for the complex reflection coefficient $\gls{sym-Gamma}$ common in electrical engineering, rather than the damping ratio and natural frequency more common in mechanical engineering. The equivalence is:
% \begin{equation}
%     \Gamma = \frac{z-1}{z+1}, \quad z = \frac{2\theta_1 + i~ \theta_4}{-2 (\theta_1+\theta_2) + i~(\theta_3+\theta_4)}
% \end{equation}
% % I am inputting zeta and omega_n corresponding to mechanical dynamcis when calculating m_sat, but I really should be inputting them corresponding to electrical dynamics. Maybe describe that?

% The gain saturation multiplier $\gls{sym-m-sat}$ is then used to modify the powertrain coefficients,
% \begin{align}
%     B_{p,sat} &= m_{sat} B_p, & K_{p,sat} &= m_{sat} K_p.
% \end{align}
% and these new control gains are then plugged into equation \eqref{eq:eom-freq-domain} to find the saturated displacement amplitude $\gls{sym-X}_{sat}$.